% Generated draft: requires mathematical and source review.
\documentclass[11pt]{article}
\usepackage[margin=25mm]{geometry}
\usepackage{fontspec}
\setmainfont{Latin Modern Roman}
\usepackage{amsmath,amssymb,mathtools,unicode-math}
\setmathfont{Latin Modern Math}
\usepackage{enumitem,needspace,xurl,hyperref}
\setlength{\parindent}{0pt}
\setlength{\parskip}{5pt}
\setlength{\emergencystretch}{3em}
\newcommand{\N}{\mathbb{N}}
\newcommand{\Z}{\mathbb{Z}}
\newcommand{\Q}{\mathbb{Q}}
\newcommand{\R}{\mathbb{R}}
\newcommand{\C}{\mathbb{C}}
\newcommand{\F}{\mathbb{F}}
\newcommand{\A}{\mathbb{A}}
\newcommand{\PP}{\mathbb{P}}
\newcommand{\E}{\mathbb{E}}
\newcommand{\eps}{\varepsilon}
\newcommand{\dd}{\,\mathrm{d}}
\begin{document}
\section*{UnsolvedMath ID: 20000869}
\textbf{Source problem code:} catalogue entry
\subsection{Definitions and mathematical statement}
This is Chapter L, ``Notes from Breakout Session II,'' a discussion among four workshop groups rather than a posed mathematical problem. The supplied aid's two mathematical results are separate proposed formalizations of parts of that exchange.

For its discrimination result, specify two simple hypotheses \(h\in\{0,1\}\) with fully known scalar response means \(f_h(x)\), common known variance \(\sigma^2>0\), and a freely repeatable input set \(\mathcal X\). At each of a fixed \(n\ge1\) trials, an input policy may depend on previous observations and independent randomization, and observes
\[
 Y_t=f_h(x_t)+\varepsilon_t,
 \qquad\varepsilon_t\text{ independent }N(0,\sigma^2).
\]
Assume
\[
 D=\max_{x\in\mathcal X}|f_1(x)-f_0(x)|<\infty
\]
is attained at \(x_*\). The aid reports that repeating \(x_*\) for all trials Blackwell-dominates every fixed-horizon adaptive policy: the latter experiment can be obtained from the former by a randomized observation transformation independent of \(h\). Its equal-prior minimum misclassification probability is
\[
 \Phi\left(-\frac{\sqrt nD}{2\sigma}\right),
\]
where \(\Phi\) is the standard normal distribution function. This includes error \(1/2\) when \(D=0\). Attainment, freely repeatable inputs, two fully specified models, known homoscedastic independent Gaussian errors, and the fixed trial budget are essential declared assumptions. The aid says a two-input composite-model counterexample disproves an unrestricted extension, but does not supply that example's parameters.

For its recording-overlap result, let a bipartite graph have cell vertices and epoch vertices, with an edge \((c,e)\) when cell \(c\) is recorded in epoch \(e\). Use the additive mean model
\[
 \E Y_{ce}=\alpha_c+\beta_e,
\]
where \(\alpha_c\) is a cell response and \(\beta_e\) is arbitrary shared epoch drift. Identifiability is structural identifiability from these observation means, with fixed centered measurement-noise laws if noise is included. For distinct observed cells \(c,c'\), the aid reports that \(\alpha_c-\alpha_{c'}\) is identifiable exactly when the two cell vertices lie in the same graph component. Componentwise transformations \(\alpha_c\mapsto\alpha_c+a\), \(\beta_e\mapsto\beta_e-a\) explain the unresolved offsets. Purely sequential epochs containing only one cell each give no between-cell contrasts under unrestricted epoch drift, even if a cell is observed repeatedly. This is not a general impossibility result when drift has additional structure.

The original discussion covers more than these models. Psychologists discuss psychophysical support for primate experiments, hiring theorists, and raw-data sharing: Alessandra argues for public availability after publication, while Judith notes possible further analyses and the wish to publish data. These are participants' views, not a legal assertion that publication places data in the public domain. Theory group I affirms associative spiking models, refers to an upcoming presentation by Fritz, requests experimental training and data, and, through Russ, urges shared-input model comparisons and documentation of failures. It says falsifiable predictions already exist and that theorists' work exceeds making testable hypotheses.

Anatomists describe a shift toward quantitative data, ask how computing and computer vision can help time-consuming tasks such as synapse counting, and discuss the effort of multielectrode recording, simultaneous activation, multichannel systems, and natural stimuli. Theory group II presents Bayesian modeling as a framework, suggests implementing priors through lateral connectivity, synaptic weights, or top-down predictions, and refers to Rajesh Rao's work for tests without specifying a paper. It suggests diffuse topographic feedback may support spatial priming and patchy feedback feature priming, explicitly acknowledging uncertainty. Tony discusses natural stimuli for adaptive models and tighter experimental predictions for hardwired models. The truncated Bayesian sentence and page number 13 do not supply missing equations or additional hypotheses. No unified theorem or completed empirical program is recovered from these notes.
\subsection{Short English statement}
This workshop discussion concerns model comparison, data sharing, quantitative anatomy, and feedback. The supplied aid separately proposes a narrowly scoped optimal Gaussian experiment and an exact graph criterion for comparing cells in the presence of arbitrary epoch drift; neither turns the full discussion into an open theorem.
\subsection{Sources}
\begin{itemize}
\item Local UnsolvedMath registry, permanent ID \(20000869\), source code \texttt{catalogue entry}; \texttt{aim-biology-notes.json}, \texttt{notes[33]}, tag \texttt{section}.
\item AIM workshop notes, \emph{Inference and Prediction in Neocortical Circuits}, version dated 24 October 2003, Chapter L, ``Notes from Breakout Session II'': \url{https://aimath.org/WWN/brain/brain.pdf}. Statement status \texttt{exact}. No PDF was accessed for this draft.
\item Discussion attributions supplied as Alessandra, Judith, Fritz, Russ, Tony, and Rajesh Rao; no fuller identities or specific reference to Rao's work is supplied. Blackwell comparison is invoked by the aid without a full bibliographic citation. No original remarks or literature field are supplied.
\item Research attempt 1: \texttt{context\_only}, problem status at run \texttt{not\_a\_problem}, rigor label \texttt{proved}, novelty confidence low; registry status \texttt{open}, classification \texttt{OPEN-TRIAGE}, suggested difficulty L2. Created and updated 14 August 2026. This machine-generated aid is not peer reviewed and must be independently verified before citation.
\item Ulam AI / UnsolvedMath Contributors; CC BY 4.0.
\end{itemize}
\end{document}
